\documentclass[12pt]{article}
\usepackage[T1]{fontenc}
\usepackage[margin=2cm]{geometry}
\pagestyle{empty}
\setlength{\parindent}{0pt}
\setlength{\parskip}{5pt}

\begin{document}

\begin{center}
{\LARGE\bfseries Savings Accounts: Answer Key}\\[2mm]
{\normalsize From the worked solutions}
\end{center}

\vspace{4mm}

\textbf{1. Amara --- Metroline Bank: 10\% simple interest, \pounds 120 for 4 years}
\begin{itemize}
  \item Interest each year: \pounds 12
  \item Interest altogether: \pounds 48
  \item Total at the end: \pounds 168
\end{itemize}

\textbf{2. Ben --- Chace: 4\% compound interest, \pounds 200 for 2 years}
\begin{itemize}
  \item Multiply by: 1.04
  \item After 1 year: \pounds 208
  \item After 2 years: \pounds 216.32
\end{itemize}

\textbf{3. Priya --- Starlight Bank: 6\% compound interest, \pounds 500}
\begin{itemize}
  \item Multiply by: 1.06
  \item After 10 years: \pounds 895.42
  \item After how many whole years is there more than \pounds 1000? \textbf{12 years}
\end{itemize}

\textbf{4. Erin --- \pounds 600 for 2 years: Virgo 8\% compound or Tesbury 6\% simple?}
\begin{itemize}
  \item Virgo pays her more.
  \item Difference: \pounds 27.84 more.
\end{itemize}

\textbf{5. Aisha --- Halifield Half-Year Saver: 2\% compound every 6 months, \pounds 500 for 1 year}
\begin{itemize}
  \item After 6 months: \pounds 510
  \item After 12 months: \pounds 520.20
  \item Is 2\% twice a year the same as 4\% once a year? \textbf{No}: it is equivalent to 4.04\% a year.
\end{itemize}

\textbf{6. Sofia --- Nationworth Loyalty Saver: 6\% simple interest with \pounds 2 monthly fee, \pounds 250 for 2 years}
\begin{itemize}
  \item Interest in 2 years: \pounds 30
  \item Fees cost: \pounds 48
  \item She is left with: \pounds 232, so she ends with less than she started with.
\end{itemize}

\end{document}